% ------------------------------------------------------------------------------------
%                               TRASMISSIONE NUMERICA
% ------------------------------------------------------------------------------------
\chapter{Trasmissione numerica}

Ricordiamo che avevamo accennato al fatto che un \textbf{sistema di telecomunicazioni}
può essere schematizzato come mostrato in figura~\ref{fig:08_01_SchemaSistemaTelecom}.

\begin{figure}[!ht]
    \centering
    \begin{tikzpicture}[
            >=Stealth, 
            thick,
            scale = 0.67,
            block/.style={
                draw=nord0, 
                fill=nord14!20, 
                rectangle, 
                minimum height=1.2cm, 
                minimum width=2.2cm, 
                align=center, 
                font=\small,
                text=nord0
            },
            sum/.style={
                draw=nord0, 
                circle, 
                inner sep=0pt, 
                minimum size=0.6cm, 
                fill=white
            },
            % Stile per i riquadri tratteggiati (Tx e Rx)
            container/.style={
                draw=nord3, 
                dashed, 
                thick, 
                inner sep=8pt
            }
            ]

            % --- Nodi Principali ---
            \node[block] (mod) {modulatore\\numerico};
            \node[block, right=1.4cm of mod] (canale) {canale};
            \node[sum, right=1.4cm of canale] (somma) {};
            \node[block, right=1.4cm of somma] (demod) {demodulatore\\numerico};

            % --- Riquadri Tratteggiati per Tx e Rx ---
            \node[container, fit=(mod), 
                label={[font=\bfseries\scriptsize, text=nord3]above:Trasmettitore (Tx)}] (tx_box) {};
            \node[container, fit=(demod), 
                label={[font=\bfseries\scriptsize, text=nord3]above:Ricevitore (Rx)}] (rx_box) {};

            % --- Etichette e Frecce ---
            \node at (somma) {\(+\)}; 

            \draw[->, nord0] ([xshift=-2cm]tx_box.west) -- (tx_box.west) 
                node[midway, above, text=nord0] {simboli};
            \draw[->, nord0] (tx_box) -- (canale);
            \draw[->, nord0] (canale) -- (somma);
            \draw[->, nord0] (somma) -- (rx_box);
            \draw[->, nord0] (rx_box.east) -- ([xshift=2cm]rx_box.east) 
                node[midway, above, text=nord0] {simboli};

            % Rumore
            \draw[->, nord0] ([yshift=-1cm]somma.south) node[below, text=nord0] {\(w(t)\)} -- (somma.south);
        \end{tikzpicture}
    \caption{Schetizzazione di un sistema di telecomunicazioni}
    \label{fig:08_01_SchemaSistemaTelecom}
\end{figure}

Facciamo alcune osservazioni su ciascun componente del sistema:

\begin{remark}[OSSERVAZIONI SULLO SCHEMA]
    Lo schema in figura mette in evidenza il fatto che:
    \begin{itemize}
        \item Lo \textbf{scopo} di un sistema di telecomunicazioni è quello di
            \textbf{trasmettere dei messaggi}, che possono essere modellati come
            delle \textbf{sequenze di simboli}.
        \item La trasmissione avviene attraverso un \textbf{canale di trasmissione}
            (e.g. cavi in rame, cavi in fibra ottica, aria ...), il quale è
            tuttavia una \textbf{fonte di disturbi} ($w(t)$).
        \item Il segnale deve essere \textit{"convertito"} in una \textbf{forma tale}
            da \textbf{poter essere trasmesso} attraverso il \textbf{canale}.
            Questa operazione è svolta da un \textbf{modulatore} (parte del blocco
            \textit{trasmettitore}), mentre in \textit{fase di ricezione} un 
            \textbf{demodulatore} (parte del blocco \textit{ricevitore}) effettua
            la \textit{trasformazione inversa} per \textit{"riconvertire"}
            il segnale nella \textit{forma desiderata}.
    \end{itemize}
\end{remark}

Per quanto riguarda i \textbf{simboli di ingresso} $a_k$, questi possono 
essere modellati come delle \textbf{variabili aleatorie discrete} 
che assumono i valori in:
\[
    \mathcal{A} = \{1,\ldots , M\}  
\]

\begin{definition}[ALFABETO]
    L'\textbf{insieme di valori} $\mathcal{A}$ che i simboli possono assumere
    è detto \textbf{alfabeto}.
\end{definition}

In particolare, possiamo immaginare che la \textbf{sorgente emetta dei simboli}
ad \textit{intervalli regolari} con \textbf{periodo di simbolo} $T_s$ (e frequenza $R_s$, detta
\textbf{symbol rate}): la sequenza di ingresso è pertanto assimilabile ad 
un \textbf{processo aleatorio discreto}.

\begin{definition}[PERIODO DI SIMBOLO]
    Il \textbf{periodo di simbolo} è l'\textit{intervallo di tempo} che
    intercorre tra \textbf{due emessioni} di simboli \textit{consecutive}.
\end{definition}

\section{Modulatore}

Concentriamoci ora sul \textbf{trasmettitore} e, in particolare, sul
\textbf{modulatore}.

Lo scopo del \textbf{modulatore} è \textbf{associare una forma d'onda} 
ad \textbf{ogni simbolo di ingresso} in modo che possa \textbf{essere compatibile}
per la \textbf{trasmisisone} attraverso il \textbf{canale}:
\[
    a_k = n\in \mathcal{A} \longrightarrow s_n(t)
\]

\begin{definition}[MODULAZIONE M-ARIA]
    Il modulatore impiega un insieme di $M$ \textbf{forme d'onda}.
    \begin{itemize}
        \item $M=2$: si parla di \textbf{modulazione binaria}.
        \item $M>2$: si parla di \textbf{modulazione M-aria}.
    \end{itemize}
\end{definition}

\subsection{Mappatura dei simboli}

Prima di poter \textit{generare le forme d'onda da trasmettere}, è spesso 
necessario effettuare una \textbf{mappatura} tra i \textbf{simboli emessi dalla sorgente}
e quelli da \textbf{trasmettere attraverso il canale}.

Per capire meglio cosa intendiamo, consideriamo il caso della \textbf{trasmissione digitale}.
Analogamente a quanto detto sopra, la sorgente emette \textit{valori binari} 
con \textbf{periodo di bit} $T_b$ (quindi con frequenza $R_b = 1 /T_b$, detta 
comunemente \textbf{bit rate}). 

\begin{remark}[TRASMISSIONE DIGITALE]
    Anzichè \textbf{inviare i singoli bit}, questi vengono
    \textbf{raggruppati} in \textbf{blocchi} di $L$ bit ciascuno
    \textbf{prima} di \textbf{passare attraverso} al \textbf{modulatore}.
\end{remark}

Con $L$ bit abbiamo $2^L$ \textbf{possibili simboli}, per cui necessiteremo di
\[
    M = 2^L
\]
\textbf{forme d'onda}. 

Vale la pena considerare anche il ragionamento inverso: se abbiamo un 
\textbf{modulatore} in grado di produrre $M$ \textbf{forme diverse}, allora 
significa che \textbf{ciascun blocco} dovrà essere \textbf{costituito} da 
$L= \log_{2}M$ \textbf{bit}.

\begin{remark}[PERIODO DI SIMBOLO E DI BIT]
    Siccome un \textbf{simbolo trasporta $L$ bit}, dovrà essere:
    \[
        T_s = L\cdot T_b
    \] 
    Ma $L=\log_{2}M$, per cui abbiamo $T_s = \log_{2}(M) \cdot T_b$.
    
    Quindi risulta:
    \[
        R_b = \frac{1}{T_b} = \frac{\log_{2}M}{T_s} = R_s\cdot \log_{2}M
    \] 
\end{remark}


Una volta \textit{stabilita la corrispondenza}, il \textbf{modulatore} si
occuperà di \textbf{associare} una \textbf{forma d'onda} a \textbf{ciascun simbolo}.

Alla luce di ciò, possiamo \textit{schematizzare il trasmettitore} come mostrato
in figura~\ref{fig:08_02_TxMapMod}.

\begin{figure}[!ht]
    \centering
    \begin{tikzpicture}[
        >=Stealth, % Punta delle frecce
        thick,     % Spessore delle linee
        % Stile per i blocchi interni
        block/.style={
            draw=nord0, 
            fill=nord14!20, % Verde Nord chiarissimo
            rectangle, 
            minimum height=1.2cm, 
            minimum width=2.2cm, 
            align=center, 
            font=\small,
            text=nord0
        },
        % Stile per il contenitore esterno
        container/.style={
            draw=nord3, % Grigio Nord scuro
            dashed, 
            thick, 
            inner sep=12pt
        }
        ]

        % --- Blocchi interni ---
        \node[block] (map) {MAP};
        \node[block, right=1.5cm of map] (mod) {Modulatore};

        % --- Contenitore esterno (Tx) ---
        \node[container, fit=(map) (mod), 
            label={[font=\bfseries\small, text=nord0]above:Trasmettitore (Tx)}] (tx) {};

        % --- Frecce e segnali ---

        % Ingresso
        \draw[->, nord0] (-2.5,0) -- (map.west) node[left=1, above] {\(a_k\)};

        % Segnale intermedio (dal MAP al Modulatore)
        \draw[->, nord0] (map) -- (mod) node[midway, above] {\(s_k\)};

        % Uscita
        \draw[->, nord0] (mod.east) -- (6,0) node[right, above] {\(S_{Tx}(t)\)};
    \end{tikzpicture}
    \caption{Schematizzazione del blocco trasmettitore}
    \label{fig:08_02_TxMapMod}
\end{figure}

\begin{remark}["PIPELINE" TX]
    A partire da un \textbf{blocco} di $L$ \textbf{bit} $a_k$ un
    \textbf{mappatore} produce un \textbf{simbolo} $s_k$ a
    cui il \textbf{modulatore} associa la \textbf{forma d'onda da trasmettere}.
\end{remark}

\subsection{Esempio}

Per capire meglio quanto spiegato sopra, consideriamo un esempio.
Immaginiamo di avere una \textbf{fonte di bit} ed un \textbf{modulatore}
con $M=4$ forme d'onda disponibili.
Per cui:
\begin{itemize}
    \item Abbiamo \textbf{4 simboli possibili}: $\mathcal{A}= \{1,2,3,4\}$.
    \item Ciascun simbolo codificherà $L=\log_{2}4 = 2$ \textbf{bit}.
\end{itemize}

La \textbf{corrispondenza} da stabilire tra simboli e bit
può quindi essere (per esempio):
\[
    \begin{aligned}
        00 &\to 1 = s_{1}(t) \\ 
        01 &\to 2 = s_{2}(t)\\ 
        10 &\to 3 = s_{3}(t)\\ 
        11 &\to 4 = s_{4}(t)
    \end{aligned}
\]
Delle possibili \textbf{forme d'onda} da \textit{associare ai simboli} sono
mostrate in figura~\ref{fig:08_03_FormeOndaM4} (in figura usiamo $T$ per indicare $T_s$).

\begin{figure}[!ht]
    \centering
    \pgfplotsset{
        axisStyle/.style={
            width=6.5cm,
            height=4.5cm,
            axis lines=middle,
            xmin=-0.2, xmax=1.2,
            ymin=-1.5, ymax=1.5,
            xtick={0, 0.5, 1},
            xticklabels={\(0\), \(T/2\), \(T\)},
            ytick={-1, 0, 1},
            yticklabels=\empty,
            xlabel={\(t\)},
            xlabel style={at={(1.1,0.5)}, anchor=west},
            ylabel style={at={(0.5,1)}, anchor=south},
            grid=major,
            grid style={dashed, gray!20},
            axis line style={-Stealth, thick},
        }
    }
    \begin{tikzpicture}
        % --- s1(t): Impulso positivo da 0 a T/2 ---
        \begin{axis}[
            name=plot1,
            at={(0, 0)},
            axisStyle,
            ylabel={\(s_1(t)\)},
            ]
            \addplot[very thick, nord10] coordinates {(-0.2,0) (0,0) (0,1) (0.5,1) (0.5,0) (1.2,0)};
            \draw[dashed, nord3] (axis cs:0,0) -- (axis cs:0,1);
        \end{axis}

        % --- s2(t): Impulso positivo da T/2 a T ---
        \begin{axis}[
            name=plot2,
            at={(plot1.east)},
            anchor=west,
            xshift=1cm,
            axisStyle,
            ylabel={\(s_2(t)\)},
            ]
            \addplot[very thick, nord10] coordinates {(-0.2,0) (0.5,0) (0.5,1) (1,1) (1,0) (1.2,0)};
            \draw[dashed, nord3] (axis cs:0.5,0) -- (axis cs:0.5,1);
        \end{axis}

        % --- s3(t): Impulso negativo da 0 a T/2 ---
        \begin{axis}[
            name=plot3,
            at={(plot1.below south west)},
            anchor=north west,
            yshift=-1cm,
            axisStyle,
            ylabel={\(s_3(t)\)},
            ]
            \addplot[very thick, nord11] coordinates {(-0.2,0) (0,0) (0,-1) (0.5,-1) (0.5,0) (1.2,0)};
            \draw[dashed, nord3] (axis cs:0,0) -- (axis cs:0,-1);
        \end{axis}

        % --- s4(t): Impulso negativo da T/2 a T ---
        \begin{axis}[
            name=plot4,
            at={(plot2.below south west)},
            anchor=north west,
            yshift=-1cm,
            axisStyle,
            ylabel={\(s_4(t)\)},
            ]
            \addplot[very thick, nord11] coordinates {(-0.2,0) (0.5,0) (0.5,-1) (1,-1) (1,0) (1.2,0)};
            \draw[dashed, nord3] (axis cs:0.5,0) -- (axis cs:0.5,-1);
        \end{axis}

    \end{tikzpicture}
    \caption{Possibili forme d'onda per il caso d'esempio}
    \label{fig:08_03_FormeOndaM4}
\end{figure}

Immaginiamo ora che la sorgente emetta la seguente \textbf{sequenza di bit}:
\[
    010010110101
\]

Il \textbf{mappatore} li raggrupperà a \textbf{gruppi di due} per 
effettuare l'\textbf{associazione dei simboli}:
\[
    \begin{matrix}
        [01] & [00] & [10] & [11] & [01] & [01] \\ 
        \downarrow &\downarrow &\downarrow &\downarrow &\downarrow &\downarrow \\ 
        2 & 1 & 3 & 4 & 2 & 2 \\ 
        \downarrow &\downarrow &\downarrow &\downarrow &\downarrow &\downarrow \\ 
        s_{2}(t) & s_{1}(t) & s_{3}(t) & s_{4}(t) & s_{2}(t) & s_{2}(t)
    \end{matrix}
\]

E il \textbf{segnale trasmesso} sarà quindi la \textit{sequenza di forme d'onda}
mostrata in figura~\ref{fig:08_04_EsempioSegnaleModulato}.

\begin{figure}[!ht]
    \centering
    \begin{tikzpicture}
        \begin{axis}[
            width=0.95\textwidth,
            height=6cm,
            axis lines=middle,
            xlabel={\(t\)},
            ylabel={\(s_{Tx}(t)\)},
            xmin=-0.5, xmax=12.5, % Asse x da 0 a 6T (dove T=2)
            ymin=-1.5, ymax=1.5,
            xtick={0, 2, 4, 6, 8, 10, 12},
            xticklabels={\(0\), \(T\), \(2T\), \(3T\), \(4T\), \(5T\), \(6T\)},
            ytick={-1, 0, 1},
            yticklabels={\(-A\), \(0\), \(A\)}, % Etichette simboliche per l'ampiezza
            xlabel style={at={(1,0.5)}, anchor=west},
            ylabel style={at={(0.5,1)}, anchor=south},
            grid=major,
            grid style={dashed, gray!20},
            axis line style={-Stealth, thick},
            ]

            % --- Tracciamento del segnale continuo s_Tx(t) ---
            \addplot[very thick, nord10] coordinates {
                    (0,0) (1,0) (1,1) (3,1) (3,0) (4,0) 
                    (4,-1) (5,-1) (5,0) (7,0) (7,-1) (8,-1) 
                    (8,0) (9,0) (9,1) (10,1) (10,0) (11,0) 
                    (11,1) (12,1) (12,0)
                };

            % Linee verticali tratteggiate per separare i simboli (ogni T=2)
            \pgfplotsinvokeforeach{2, 4, 6, 8, 10}{
                \draw[dashed, gray!50] (axis cs:#1,-1.5) -- (axis cs:#1,1.5);
            }

            % Etichette dei bit raggruppati (sopra)
            \node[above, font=\small] at (axis cs:1,1.2) {01};
            \node[above, font=\small] at (axis cs:3,1.2) {00};
            \node[above, font=\small] at (axis cs:5,1.2) {10};
            \node[above, font=\small] at (axis cs:7,1.2) {11};
            \node[above, font=\small] at (axis cs:9,1.2) {01};
            \node[above, font=\small] at (axis cs:11,1.2) {01};

            % Etichette dei simboli (sotto)
            \node[below, font=\small, text=nord11] at (axis cs:1,-1.2) {2};
            \node[below, font=\small, text=nord11] at (axis cs:3,-1.2) {1};
            \node[below, font=\small, text=nord11] at (axis cs:5,-1.2) {3};
            \node[below, font=\small, text=nord11] at (axis cs:7,-1.2) {4};
            \node[below, font=\small, text=nord11] at (axis cs:9,-1.2) {2};
            \node[below, font=\small, text=nord11] at (axis cs:11,-1.2) {2};
        \end{axis}
    \end{tikzpicture}
    \caption{Segnale trasmesso}
    \label{fig:08_04_EsempioSegnaleModulato}
\end{figure}

\subsection{PWM}

Vediamo ora un esempio di modulazione di
tipo \textbf{PWM}: \textbf{Pulse Width Modulation}.

\begin{remark}[PWM]
    In questo tipo di modulazione si utilizzano \textbf{onde quadre} con
    \textbf{diverse larghezze degli impulsi}.

    E' ampiementa utilizzata in \textbf{elettronica digitale} perchè
    permette di \textit{controllare la potenza di motori elettrici}
    usando \textit{segnali digitali}.
\end{remark}

Usiamo per esempio $M=4$ con i segnali seguenti, mostrati 
anche in figura~\ref{fig:08_05_EsempioPWM}.
\[
    \begin{aligned}
        s_{1}(t) &= \begin{cases}
            A & 0 < t <\frac{T}{4} \\ 
            0 & \text{altrimenti}
        \end{cases} \\ 
        s_{2}(t) &= \begin{cases}
            A & 0 < t <\frac{T}{2} \\ 
            0 & \text{altrimenti}
        \end{cases} \\ 
        s_{3}(t) &= \begin{cases}
            A & 0 < t <\frac{3T}{4} \\ 
            0 & \text{altrimenti}
        \end{cases} \\ 
        s_{4}(t) &= \begin{cases}
            A & 0 < t < T \\ 
            0 & \text{altrimenti}
        \end{cases} \\ 
    \end{aligned}
\]

\begin{figure}[!ht]
    \centering
    \pgfplotsset{
        pwmStyle/.style={
            width=6.5cm,
            height=3.5cm,
            axis lines=middle,
            xmin=-0.05, xmax=1.15,
            ymin=-0.1, ymax=1.3,
            xtick={0, 0.25, 0.5, 0.75, 1},
            xticklabels={\(0\), \(\frac{T}{4}\), \(\frac{T}{2}\), \(\frac{3T}{4}\), \(T\)},
            ytick={0, 1},
            yticklabels={\(0\), \(A\)},
            ylabel style={at={(0,1)}, anchor=south east},
            grid=major,
            grid style={dashed, gray!20},
            axis line style={-Stealth, thick},
        }
    }

    \begin{tikzpicture}

        % --- s1(t) ---
        \begin{axis}[name=ax1, pwmStyle, ylabel={\(s_1(t)\)}]
            \addplot[very thick, nord10] coordinates {(0,0) (0,1) (0.25,1) (0.25,0) (1,0)};
        \end{axis}

        % --- s2(t) ---
        \begin{axis}[name=ax2, pwmStyle, at={(ax1.east)}, anchor=west, xshift=1cm, ylabel={\(s_2(t)\)}]
            \addplot[very thick, nord10] coordinates {(0,0) (0,1) (0.5,1) (0.5,0) (1,0)};
        \end{axis}

        % --- s3(t) ---
        \begin{axis}[name=ax3, pwmStyle, at={(ax1.south west)}, 
            anchor=north west, yshift=-1cm, ylabel={\(s_3(t)\)}]
            \addplot[very thick, nord10] coordinates {(0,0) (0,1) (0.75,1) (0.75,0) (1,0)};
        \end{axis}

        % --- s4(t) ---
        \begin{axis}[name=ax4, pwmStyle, at={(ax2.south west)}, 
            anchor=north west, yshift=-1cm, ylabel={\(s_4(t)\)}]
            \addplot[very thick, nord10] coordinates {(0,0) (0,1) (1,1) (1,0)};
        \end{axis}
    \end{tikzpicture}
    \caption{Esempio di forme d'onda per modulazione PWM}
    \label{fig:08_05_EsempioPWM}
\end{figure}

Supponiamo di voler \textit{ricavare una base ortonormale} per questi segnali.
Notiamo subito che:

\begin{remark}
    I segnali sono tutti \textbf{linearmente indipendenti} (non siamo 
    in grado di scrivere uno come combinazione lineare degli altri), 
    tuttavia \textbf{non sono ortogonali}.
\end{remark}

Per ottenere una \textbf{base ortonormale} potremmo usare il
procedimento di Gram-Schmidt, tuttavia risulta spesso particolarmente
laborioso. Possiamo invece \textbf{scegliere intelligentemente} una base
del tipo:
\[
    \begin{aligned}
        \phi_{1}(t) &= \begin{cases}
            A' & 0 < t <\frac{T}{4} \\ 
            0 & \text{altrimenti}
        \end{cases} \\ 
        \phi_{2}(t) &= \begin{cases}
            A' & \frac{T}{4} < t <\frac{T}{2} \\ 
            0 & \text{altrimenti}
        \end{cases} \\ 
        \phi_{3}(t) &= \begin{cases}
            A' & \frac{T}{2} < t <\frac{3T}{4} \\ 
            0 & \text{altrimenti}
        \end{cases} \\ 
        \phi_{4}(t) &= \begin{cases}
            A' & \frac{3T}{4} < t < T \\ 
            0 & \text{altrimenti}
        \end{cases} \\ 
    \end{aligned}
\]

I segnali di questa base sono mostrati in figura~\ref{fig:08_06_EsempioPWMBaseOrto}.

\begin{figure}[!ht]
    \centering
    \pgfplotsset{
        basisStyle/.style={
            width=6.5cm,
            height=3.5cm,
            axis lines=middle,
            xmin=-0.05, xmax=1.05,
            ymin=-0.1, ymax=1.3,
            xtick={0, 0.25, 0.5, 0.75, 1},
            xticklabels={\(0\), \(\frac{T}{4}\), \(\frac{T}{2}\), \(\frac{3T}{4}\), \(T\)},
            ytick={0, 1},
            yticklabels={\(0\), \(A'\)},
            ylabel style={at={(0,1)}, anchor=south east},
            grid=major,
            grid style={dashed, gray!20},
            axis line style={-Stealth, thick},
        }
    }

    \begin{tikzpicture}

        % --- phi_1(t) ---
        \begin{axis}[name=ax5, basisStyle, ylabel={\(\phi_1(t)\)}]
            \addplot[very thick, nord14] coordinates {(0,0) (0,1) (0.25,1) (0.25,0) (1,0)};
        \end{axis}

        % --- phi_2(t) ---
        \begin{axis}[name=ax6, basisStyle, at={(ax5.east)}, anchor=west, xshift=1cm, ylabel={\(\phi_2(t)\)}]
            \addplot[very thick, nord14] coordinates {(0,0) (0.25,0) (0.25,1) (0.5,1) (0.5,0) (1,0)};
        \end{axis}

        % --- phi_3(t) ---
        \begin{axis}[name=ax7, basisStyle, at={(ax5.south west)}, 
            anchor=north west, yshift=-1cm, ylabel={\(\phi_3(t)\)}]
            \addplot[very thick, nord14] coordinates {(0,0) (0.5,0) (0.5,1) (0.75,1) (0.75,0) (1,0)};
        \end{axis}

        % --- phi_4(t) ---
        \begin{axis}[name=ax8, basisStyle, at={(ax6.south west)}, 
            anchor=north west, yshift=-1cm, ylabel={\(\phi_4(t)\)}]
            \addplot[very thick, nord14] coordinates {(0,0) (0.75,0) (0.75,1) (1,1) (1,0)};
        \end{axis}

    \end{tikzpicture}
    \caption{Base scelta per la PWM d'esempio}
    \label{fig:08_06_EsempioPWMBaseOrto}
\end{figure}

Infine notiamo che i segnali sono tutti a \textbf{supporto disgiunto}, per cui il loro
\textbf{prodotto interno} è \textbf{sicuramente nullo} e quindi \textbf{formano una base ortogonale}. 
La scelta è stata intelligente perchè abbiamo banalmente:
\[
    \begin{aligned}
        s_{1}(t) &= \alpha \phi_{1}(t) \\ 
        s_{2}(t) &= \alpha \left[ \phi_{1}(t) + \phi_{2}(t) \right] \\ 
        s_{3}(t) &= \alpha \left[ \phi_{1}(t) + \phi_{2}(t) + \phi_{3}(t) \right] \\ 
        s_{4}(t) &= \alpha \left[ \phi_{1}(t) + \phi_{2}(t) + \phi_{3}(t) + \phi_{4}(t) \right] \\ 
    \end{aligned}
\]

Ci resta infine solo da \textbf{determinare l'ampiezza} affinchè la norma dei vettori 
sia unitaria. Per tutti e quattro i segnali l'energia risulta:
\[
    E_{\phi} = A'^2 \frac{T}{4}
\]

Per cui scegliamo:
\[
    A' = \frac{2}{\sqrt{T}} = \frac{2\sqrt{T}}{T}
\]

\section{Canale}

Vediamo ora \textit{cosa succede al segnale} \textbf{durante la trasmissione}.

Abbiamo già accennato al fatto che il \textbf{canale introduce rumore nel segnale}:
se consideriamo un segnale simile a quello dell'esempio visto sopra (composto dalle
\textit{stesse forme d'onda}), possiamo allora trovarci in una situazione come quella 
illustrata in figura~\ref{fig:08_07_RumoreCanale}.

\begin{figure}[!ht]
    \centering
    \begin{tikzpicture}[
        % Definiamo le funzioni per il segnale e il rumore
        declare function={
            % Segnale pulito a tratti (1 per i positivi, -1 per i negativi, 0 altrove)
            sig(\x) = and(\x>0.5, \x<2)*1 + and(\x>2.5, \x<4)*(-1) + and(\x>5, \x<6)*1 + 
            and(\x>6.5, \x<7.5)*1;
            % Rumore simulato con una somma di sinusoidi a frequenze diverse
            noise(\x) = 0.2*sin(deg(2*pi*0.4*\x)) + 0.15*sin(deg(2*pi*1.2*\x)) 
            + 0.1*sin(deg(2*pi*2.7*\x)) + 0.05*sin(deg(2*pi*5.1*\x));
        }
        ]
        % --- Plot 1: Segnale trasmesso ---
        \begin{axis}[
            name=ax1,
            width=0.9\textwidth,
            height=3.5cm,
            axis lines=middle,
            xmin=0, xmax=8,
            ymin=-1.8, ymax=2.5,
            xtick={1,2,3,4,5,6},
            xticklabels={\(T\), \(2T\), \(3T\), \(4T\), \(5T\), \(6T\)},
            ytick=\empty,
            ylabel={\(s_{Tx}(t)\)},
            xlabel={},
            ]
            % Tracciamento del segnale ideale tramite coordinate per avere spigoli netti
            \addplot[very thick, nord0] coordinates {
                    (0,0) (0.5,0) (0.5,1) (2,1) (2,0) (2.5,0) (2.5,-1) (4,-1) (4,0) 
                    (5,0) (5,1) (6,1) (6,0) (6.5,0) (6.5,1) (7.5,1) (7.5,0) (8,0)
                };
        \end{axis}
        % Titolo allineato a sinistra
        \node[anchor=south west, font=\bfseries, text=nord0] 
            at (ax1.north west) {Segnale trasmesso};

        % --- Plot 2: Rumore ---
        \begin{axis}[
            name=ax2,
            at={(ax1.below south west)},
            anchor=north west,
            yshift=-0.5cm, % Spazio tra i grafici
            width=0.9\textwidth,
            height=3.5cm,
            axis lines=middle,
            xmin=0, xmax=8,
            ymin=-1.8, ymax=2.5,
            xtick={1,2,3,4,5,6},
            xticklabels={\(T\), \(2T\), \(3T\), \(4T\), \(5T\), \(6T\)},
            ytick=\empty,
            ylabel={\(w_{c}(t)\)},
            xlabel={},
            ]
            \addplot[very thick, nord11, domain=0:8, samples=400] {noise(x)};
        \end{axis}
        \node[anchor=south west, font=\bfseries, text=nord0] 
            at (ax2.north west) {Rumore};

        % --- Plot 3: Segnale ricevuto ---
        \begin{axis}[
            name=ax3,
            at={(ax2.below south west)},
            anchor=north west,
            yshift=-0.5cm,
            width=0.9\textwidth,
            height=3.5cm,
            axis lines=middle,
            xmin=0, xmax=8,
            ymin=-1.8, ymax=2.5,
            xtick={1,2,3,4,5,6},
            xticklabels={\(T\), \(2T\), \(3T\), \(4T\), \(5T\), \(6T\)},
            ytick=\empty,
            ylabel={\(s_{Rx}(t)\)},
            xlabel={},
            ]
            % 1. Contorno nero del segnale originale (per riferimento)
            \addplot[thick, nord0] coordinates {
                    (0,0) (0.5,0) (0.5,1) (2,1) (2,0) (2.5,0) (2.5,-1) (4,-1) (4,0) 
                    (5,0) (5,1) (6,1) (6,0) (6.5,0) (6.5,1) (7.5,1) (7.5,0) (8,0)
                };
            % 2. Segnale ricevuto (Segnale + Rumore) in verde
            \addplot[very thick, nord14, domain=0:8, samples=400] {sig(x) + noise(x)};
        \end{axis}
        \node[anchor=south west, font=\bfseries, text=nord0] 
            at (ax3.north west) {Segnale ricevuto dopo il canale};

    \end{tikzpicture}
    \caption{Introduzione di rumore nel segnale durante la trasmissione}
    \label{fig:08_07_RumoreCanale}
\end{figure}

E' facile intuire che il rumore mostrato in figura è 
assimilabile ad un \textbf{processo aleatorio continuo}.

\begin{remark}[CANALE COME FILTRO]
    In generale, è possibile che il canale, \textit{oltre ad introdurre del rumore},
    applichi \textbf{anche una trasformazione} $h_c(t)$ al segnale.
    Pensiamo banalmente alla \textbf{perdita di potenza} (\textit{riduzione dell'ampiezza}
    del segnale) nella trasmissione sulle \textbf{lunghe distanze}
\end{remark}

\subsection{Rappresentazione vettoriale}

Ciascuna \textit{forma d'onda} $s_n(t)$ ($n\in \{1,2,3,4\}$) giunge a destinazione 
\textbf{corrotta dal rumore} $w(t)$, per cui il \textbf{ricevitore} legge 
un segnale $r(t)$ dato da:
\[
    r(t) = s_n(t) + w(t)
\]

Ora vogliamo analizzare meglio la situazione e, in particolare, vogliamo studiare 
le  \textbf{rappresentazioni vettoriali} del \textbf{segnale originale}, 
del \textbf{rumore} e del \textbf{segnale ricevuto}. 

Per farlo, abbiamo bisogno di \textbf{scegliere una base ortonormale adeguata}.

\begin{remark}[SCELTA DELLA BASE]
    E' intuitivo capire che conviene scegliere una base che 
    \textbf{contenga le forme d'onda del modulatore}, visto che queste
    sono quelle che siamo interessati a ricevere.
\end{remark}

E' facile notare che i quattro segnali \textit{non sono linearmente indipendenti}, infatti:
\[
    \begin{aligned}
        s_{3}(t) &= -s_{1}(t) \\ 
        s_{4}(t) &= -s_{2}(t)
    \end{aligned}
\]
Ma notiamo anche $s_{1}(t) \perp s_{2}(t)$, infatti:
\[
    \langle \mathbf{s_{1}} , \mathbf{s_{2}} \rangle = \int_{0}^T s_{1}(t) \overline{s_{2}}(t) dt = 0
\]
Per cui la \textit{nostra scelta} di \textbf{base ortonormale} sarà:
\[
    \mathcal{B} = \{\phi_{1}, \phi_{2}\} = 
        \left\{ \frac{\mathbf{s_{1}}}{\sqrt{E_{s_{1}}}}, \frac{\mathbf{s_{2}}}{\sqrt{E_{s_{2}}}} \right\}
\]

Fatta questa scelta, possiamo \textit{rappresentare graficamente} i vettori $s_n(t)$ 
e $r(t)$, come mostrato in figura~\ref{fig:08_08_CostellazioneRumore}.

\begin{figure}[!ht]
    \centering
    \begin{tikzpicture}[
        >=Stealth,
        thick,
        scale = 0.8,
        % Stile per i simboli ideali (punti neri)
        signal/.style={circle, fill=nord0, inner sep=3pt},
        % Stile per i segnali ricevuti (punti verdi con bordo nero)
        noisy/.style={circle, draw=black, fill=nord11, inner sep=2.5pt, thick}
        ]

        % --- Assi Cartesiani ---
        \draw[->] (-4.5, 0) -- (4.5, 0) node[right] {$\phi_{1}$};
        \draw[->] (0, -4) -- (0, 4.5) node[above] {$\phi_{2}$};

        % --- Simboli Ideali (Punti Neri) ---
        \node[signal] (s1) at (3, 0) {};
        \node[signal] (s2) at (0, 3) {};
        \node[signal] (s3) at (-3, 0) {};
        \node[signal] (s4) at (0, -3) {};

        % Etichette dei simboli ideali
        \node[above=0.1cm] at (s1) {\(\mathbf{s}_1\)};
        \node[right=0.1cm] at (s2) {\(\mathbf{s}_2\)};
        \node[above right=0.05cm and 0.1cm] at (s3) {\(\mathbf{s}_3\)};
        \node[right=0.1cm] at (s4) {\(\mathbf{s}_4\)};

        % --- Segnali Ricevuti (Punti Verdi) ---
        % Punto vicino a s1
        \node[noisy] at (2.2, 0.4) {};

        % Due punti vicino a s2
        \node[noisy] at (-0.4, 3.5) {};
        \node[noisy] at (0.5, 2.4) {};

        % Punto vicino a s3
        \node[noisy] at (-3.6, 0.5) {};

        % Punto vicino a s4
        \node[noisy] at (0.5, -2.4) {};

    \end{tikzpicture}
    \caption{Segnale trasmesso e ricevuto nella rappresentazione vettoriale}
    \label{fig:08_08_CostellazioneRumore}
\end{figure}

\begin{remark}[COSTELLAZIONE: EFFETTO DEL RUMORE]
    A causa del rumore, le \textbf{forme d'onda ricevute} si
    \textbf{discostano} dai \textbf{valori "puri"} $\mathbf{s}_n$.
\end{remark}

Ne consegue che il \textbf{ricevitore} dovrà adottare delle \textbf{strategie}
per \textbf{determinare la forma d'onda corretta} da associare
a ciascuna \textbf{osservazione} $r(t)$.

Si tratta cioè di formulare $M$ (nel nostro caso $4$) \textbf{ipotesi}:
\[
    H_n : r(t) = s_n(t) + w(t) \text{ , } n = 1,2,\ldots ,M
\]
Che in \textit{forma vettoriale} diventano:
\[
    H_n : \mathbf{r} = \mathbf{s_n} + \mathbf{w} \text{ , } n = 1,2,\ldots ,M
\]

E determinare \textbf{quale è la più plausibile}.
Vedremo più avanti come implementare tali strategie.

\subsection{Caratterizzazione del rumore e canale AWGN}

Abbiamo già detto che il rumore è un \textbf{processo aleatorio}, ma
aggiungiamo ora alcune \textbf{assunzioni} che ci permettano di 
\textbf{caratterizzare} meglio il rumore.

\begin{remark}[ASSUNZIONI SUL RUMORE]
    E' ragionevole assumere che:
    \begin{enumerate}
        \item Il rumore \textbf{contenga tutte le frequenze} possibili 
            (cioè ha banda \textit{illimitata} $[0,+\infty)$).
        \item Il rumore abbia la \textbf{stessa potenza} su \textbf{ogni frequenza} 
            (cioè \textit{non c'è una frequenza predominante}).
        \item Fissato un istante di tempo $t$ arbitrario, il rumore sia 
            \textbf{distribuito normalmente} con \textbf{media nulla}.
        \item Dati \textbf{due istanti} di tempo $t_{1}$ e $t_{2}$, i
            valori $w(t_{1})$ e $w(t_{2})$ siano \textbf{indipendenti}.
    \end{enumerate}
\end{remark}

Sotto queste ipotesi, il rumore è detto \textbf{WGN}: \textbf{White Gaussian Noise}.
Inoltre, essendo il rumore \textbf{additivo}, nel senso che 
\textit{si somma al segnale} (\textit{non} introduce \textit{distorsioni moltiplicative}),
si parla di \textbf{Additive White Gaussian Noise}.

\begin{definition}[CANALE AWGN]
    Si parla di \textbf{modello} di \textbf{canale AWGN} quando
    si assumono vere le ipotesi fatte sopra.
\end{definition}

Abbiamo assunto che la \textbf{media sia nulla}: è infatti ragionevole 
ipotizzare che \textit{variazioni additive} e \textit{variazioni sottrattive} 
siano \textit{egualmente probabili}.
Per concludere la caratterizzazione dobbiamo quindi solo dare una \textbf{stima} della
\textbf{varianza del rumore}. In generale, si pone:

\begin{definition}[VARIANZA DEL RUMORE]
    \[
        \sigma^2_{w} = \frac{N_{0}}{2}
    \] 
    Dove $N_{0}$ è la \textbf{densità spettrale di potenza} del rumore.
\end{definition}

In generale, $N_{0}$ \textit{dipende dal tipo di rumore}. Per il rumore termico, per
esempio, abbiamo $N_{0} = k\cdot T$, dove $k$ è la \textit{costante di Boltzmann} e 
$T$ è la \textit{temperatura} del mezzo in Kelvin.

\subsection{Signal to Noise Ratio}

Per \textbf{quantificare l'effetto del rumore} sul \textbf{segnale ricevuto}
si definisce il \textbf{rapporto segnale-rumore}:

\begin{definition}[SIGNAL TO NOISE RATIO - SNR]
    Dati un \textit{segnale} $s(t)$ e del \textit{rumore} $n(t)$, si
    definisce il \textbf{rapporto segnale-rumore} come il \textbf{rapporto} tra
    la \textbf{potenza del segnale} e la \textbf{potenza del rumore}:
    \[
        \text{SNR} \coloneqq \frac{P_s}{P_n}
    \] 
\end{definition}

Dato che noi considereremo \textbf{canali AWGN}, è utile tenere a mente che:

\begin{prop}[POTENZA RUMORE BIANCO]
    La \textbf{potenza media} di un \textbf{rumore gaussiano} a
    \textbf{media nulla} è:
    \[
        P_s = P_{w} = \sigma^2_{w}
    \] 
\end{prop}

In ogni caso, è chiaro che \textbf{maggiore} è il valore di \textbf{SNR} e
\textbf{più pulito è il segnale}, come mostrato in figura~\ref{fig:08_09_EsempiSNR}.

\begin{figure}[!ht]
    \centering
    \pgfplotsset{
        axisLeft/.style={
            width=6.5cm,
            height=3.5cm,
            axis lines=middle,
            xmin=0, xmax=1.1,
            ymin=-2, ymax=2,
            xtick={0, 0.2, 0.4, 0.6, 0.8, 1},
            ytick={-2, -1, 0, 1, 2},
            yticklabels=\empty,
            grid=major,
            grid style={dashed, gray!30},
            axis line style={-Stealth, thick},
        },
        axisRight/.style={
            width=6.5cm,
            height=3.5cm,
            axis lines=box,
            xmin=0, xmax=1.1,
            ymin=-2, ymax=2,
            xtick={0, 0.2, 0.4, 0.6, 0.8, 1},
            ytick={-2, -1, 0, 1, 2},
            yticklabels=\empty,
            grid=major,
            grid style={dashed, gray!30},
            axis line style={thick},
        }
    }
    \begin{tikzpicture}[
        declare function={
            % Segnale sinusoidale (4 cicli in t=[0,1])
            sig(\x) = cos(deg(2*pi*4*\x));
            % Rumore simulato (somma di sinusoidi ad alta frequenza)
            noise(\x) = 0.3*sin(deg(2*pi*43*\x)) + 0.2*sin(deg(2*pi*71*\x)) 
            + 0.15*sin(deg(2*pi*113*\x)) + 0.1*sin(deg(2*pi*151*\x)) + 0.1*sin(deg(2*pi*193*\x));
        },
        % Stile per il riquadro giallo (questo è un normale stile TikZ, va bene qui)
        snrBox/.style={
            fill=nord10!20,
            draw=black, 
            thick, 
            inner sep=4pt, 
            font=\bfseries, 
            anchor=south west
        }
        ]

        % ==========================================
        % RIGA 1: SNR = 10 dB (Rumore molto alto)
        % ==========================================
        \begin{axis}[name=ax11, at={(0,0)}, anchor=north west, axisLeft]
            \addplot[thick, nord12, domain=0:1, samples=100] {sig(x)};
            \addplot[thin, nord10, domain=0:1, samples=400] {sig(x) + 1.0*noise(x)};
        \end{axis}
        \node[snrBox] (t1) at ([xshift=1cm, yshift=0.2cm]ax11.north west) {SNR = 10dB};

        % ==========================================
        % RIGA 2: SNR = 13 dB (Rumore medio)
        % ==========================================
        \begin{axis}[name=ax21, at={(ax11.south west)}, anchor=north west, yshift=-1cm, axisLeft]
            \addplot[thick, nord12, domain=0:1, samples=100] {sig(x)};
            \addplot[thin, nord10, domain=0:1, samples=400] {sig(x) + 0.6*noise(x)};
        \end{axis}
        \node[snrBox] (t2) at ([xshift=1cm, yshift=0.2cm]ax21.north west) {SNR = 13dB};

        % ==========================================
        % RIGA 3: SNR = 20 dB (Rumore basso)
        % ==========================================
        \begin{axis}[name=ax31, at={(ax21.south west)}, anchor=north west, yshift=-1cm, axisLeft]
            \addplot[thick, nord12, domain=0:1, samples=100] {sig(x)};
            \addplot[thin, nord10, domain=0:1, samples=400] {sig(x) + 0.25*noise(x)};
        \end{axis}
        \node[snrBox] (t3) at ([xshift=1cm, yshift=0.2cm]ax31.north west) {SNR = 20dB};
    \end{tikzpicture}
    \caption{Confronto tra diversi livelli di Signal-to-Noise Ratio (SNR)}
    \label{fig:08_09_EsempiSNR}
\end{figure}

\begin{remark}
    Dato che la \textbf{potenza} del \textbf{segnale} può essere 
    \textbf{diversi ordini di grandezza} più grande del \textbf{rumore},
    il valore di SNR è \textbf{spesso espresso in decibel}.
\end{remark}
